\documentclass[a4paper]{amsproc}

\usepackage{parskip}
% Preserve article's smaller abstract text and margins, without its indent.
% AMS classes use a different abstract implementation (not quotation).
\makeatletter
\@ifclassloaded{article}{%
\patchcmd{\abstract}{\quotation}{\quote}{}{%
  \PackageError{paper}{Cannot remove abstract indentation}{}%
}
\patchcmd{\endabstract}{\endquotation}{\endquote}{}{%
  \PackageError{paper}{Cannot remove abstract indentation}{}%
}
}{}
% The Markdown metadata title precedes the abstract.  Unlike the native AMS
% title routine, titling does not flush the collected abstract box.
\@ifclassloaded{amsproc}{%
  \apptocmd{\endabstract}{\@setabstract}{}{%
    \PackageError{paper}{Cannot display the AMS abstract after the title}{}%
  }%
}{}
\makeatother

\usepackage[backend=biber,style=alphabetic]{biblatex}
\addbibresource{bibliography.bib}

\usepackage
{amsmath,amssymb,hyperref,mathpartir,mathtools,minted,
simplebnf,stmaryrd,tikz-cd,verified-badges}

\usepackage{amsthm}
% Use amsart's theorem typography while retaining the amsproc page layout.
% Unlike generic amsthm, AMS classes use half-baseline theorem spacing,
% including for remarks. Keep existing environment names and counters.
\newtheoremstyle{plain}
  {.5\baselineskip plus .2\baselineskip minus .2\baselineskip}
  {.5\baselineskip plus .2\baselineskip minus .2\baselineskip}
  {\itshape}{0pt}{\bfseries}{.}{5pt plus 1pt minus 1pt}{}
\newtheoremstyle{definition}
  {.5\baselineskip plus .2\baselineskip minus .2\baselineskip}
  {.5\baselineskip plus .2\baselineskip minus .2\baselineskip}
  {\normalfont}{0pt}{\bfseries}{.}{5pt plus 1pt minus 1pt}{}
\newtheoremstyle{remark}
  {.5\baselineskip plus .2\baselineskip minus .2\baselineskip}
  {.5\baselineskip plus .2\baselineskip minus .2\baselineskip}
  {\normalfont}{0pt}{\itshape}{.}{5pt plus 1pt minus 1pt}{}
\theoremstyle{plain}
\newtheorem{theorem}{Theorem}
\newtheorem*{theorem*}{Theorem}
\newtheorem{lemma}{Lemma}
\newtheorem{proposition}[lemma]{Proposition}
\newtheorem{corollary}{Corollary}
\theoremstyle{definition}
\newtheorem{definition}{Definition}
\theoremstyle{remark}
\newtheorem*{remark}{Remark}

\usepackage{newunicodechar}
\newunicodechar{Γ}{\ensuremath{\Gamma}}
\newunicodechar{Δ}{\ensuremath{\Delta}}
\newunicodechar{Π}{\ensuremath{\Pi}}
\newunicodechar{λ}{\ensuremath{\lambda}}
\newunicodechar{ᵀ}{\textsuperscript{T}}
\newunicodechar{ᵁ}{\textsuperscript{U}}
\newunicodechar{ᵗ}{\textsuperscript{t}}
\newunicodechar{ᵘ}{\textsuperscript{u}}
\newunicodechar{ᵣ}{\textsubscript{r}}
\newunicodechar{⊔}{\ensuremath{\sqcup}}
% Native Unicode spelling, with the original article's math glyphs/fonts.
\newunicodechar{β}{\ensuremath{\beta}}
\newunicodechar{α}{\ensuremath{\alpha}}
\newunicodechar{κ}{\ensuremath{\kappa}}
\newunicodechar{ω}{\ensuremath{\omega}}
\newunicodechar{ℵ}{\ensuremath{\aleph}}
\newunicodechar{ε}{\ensuremath{\varepsilon}}
\newunicodechar{→}{\ensuremath{\to}}
\newunicodechar{⇒}{\ensuremath{\Rightarrow}}
\newunicodechar{↔}{\ensuremath{\leftrightarrow}}
\newunicodechar{↦}{\ensuremath{\mapsto}}
\newunicodechar{⊢}{\ensuremath{\vdash}}
\newunicodechar{∈}{\ensuremath{\in}}
\newunicodechar{∉}{\ensuremath{\notin}}
\newunicodechar{∀}{\ensuremath{\forall}}
\newunicodechar{∃}{\ensuremath{\exists}}
\newunicodechar{≤}{\ensuremath{\le}}
\newunicodechar{≥}{\ensuremath{\ge}}
\newunicodechar{≠}{\ensuremath{\neq}}
\newunicodechar{×}{\ensuremath{\times}}
\newunicodechar{⊆}{\ensuremath{\subseteq}}
\newunicodechar{⊂}{\ensuremath{\subset}}
\newunicodechar{∪}{\ensuremath{\cup}}
\newunicodechar{∩}{\ensuremath{\cap}}
\newunicodechar{∧}{\ensuremath{\land}}
\newunicodechar{∨}{\ensuremath{\lor}}
\newunicodechar{¬}{\ensuremath{\neg}}
\newunicodechar{⋅}{\ensuremath{\cdot}}
\newunicodechar{≜}{\ensuremath{\triangleq}}
\newunicodechar{⇉}{\ensuremath{\rightrightarrows}}
\newunicodechar{⇔}{\ensuremath{\Leftrightarrow}}
\newunicodechar{□}{\ensuremath{\square}}
\newunicodechar{△}{\ensuremath{\triangle}}
% TeX does not automatically break compound names at a Unicode en dash.
% Keep the printed dash, allowing a line break immediately after it.
\newunicodechar{–}{\textendash\allowbreak}
\newunicodechar{“}{``}
\newunicodechar{”}{''}
\newunicodechar{∞}{\ensuremath{\infty}}
\newunicodechar{≡}{\ensuremath{\equiv}}
\newunicodechar{ℓ}{\ensuremath{\ell}}
\newunicodechar{ν}{\ensuremath{\nu}}
\newunicodechar{ρ}{\ensuremath{\rho}}
\newunicodechar{δ}{\ensuremath{\delta}}
\newunicodechar{∗}{\ensuremath{*}}
\newunicodechar{ℕ}{\ensuremath{\mathbb{N}}}
\newunicodechar{ℤ}{\ensuremath{\mathbb{Z}}}
\newunicodechar{ℚ}{\ensuremath{\mathbb{Q}}}
\newunicodechar{ℝ}{\ensuremath{\mathbb{R}}}
\newunicodechar{↝}{\ensuremath{\leadsto}}
\newunicodechar{⦇}{\ensuremath{\llparenthesis}}
\newunicodechar{⦈}{\ensuremath{\rrparenthesis}}
\newunicodechar{⟦}{\ensuremath{\llbracket}}
\newunicodechar{⟧}{\ensuremath{\rrbracket}}
\newunicodechar{⟨}{\ensuremath{\langle}}
\newunicodechar{⟩}{\ensuremath{\rangle}}
\newunicodechar{σ}{\ensuremath{\sigma}}
\newunicodechar{Ω}{\ensuremath{\Omega}}
\newunicodechar{φ}{\ensuremath{\varphi}}
\newunicodechar{⊥}{\ensuremath{\bot}}
\newunicodechar{⊤}{\ensuremath{\top}}
\newunicodechar{∘}{\ensuremath{\circ}}
\newunicodechar{∅}{\ensuremath{\varnothing}}
\newunicodechar{≔}{\ensuremath{\coloneqq}}
\newunicodechar{↑}{\ensuremath{\uparrow}}
\newunicodechar{⟶}{\ensuremath{\longrightarrow}}
\newunicodechar{⟺}{\ensuremath{\Longleftrightarrow}}
\newunicodechar{≃}{\ensuremath{\simeq}}
\newunicodechar{⊬}{\ensuremath{\nvdash}}
\newunicodechar{⋯}{\ensuremath{\cdots}}
\newunicodechar{…}{\ensuremath{\ldots}}
\newunicodechar{∏}{\ensuremath{\prod}}
\newunicodechar{⋁}{\ensuremath{\bigvee}}

\DeclareRobustCommand{\PTSrk}{\textnormal{PTS-}\ensuremath{rk}}
\pdfstringdefDisableCommands{\def\PTSrk{PTS-rk}}

% Standalone offside %tikzcd blocks already supply their own math mode.
\BeforeBeginEnvironment{tikzcd}{\begin{center}}
\AfterEndEnvironment{tikzcd}{\end{center}}

% A float already supplies external spacing and a caption gap. Use mathpartir's
% paragraph layout inside figures, not a second $$ display with an empty line
% and display skips above/below the rules. This binding is local to each float.
\newenvironment{figuremathpar}[1][]{%
  \par\noindent\vbox\bgroup
  \parskip=0pt\hsize=\linewidth\centering
  \MathparLineskip #1\noindent$\displaystyle\MathparBindings
}{\unskip$\par\egroup\par}
\AtBeginEnvironment{figure}{%
  \let\mathpar\figuremathpar
  \let\endmathpar\endfiguremathpar
}

\newcommand{\rk}[1]{\mathsf{rk}_{#1}}
\newcommand{\PTS}{\mathsf{PTS}_{rk}}
\newcommand{\Rt}{\widetilde{\mathsf R}}
\newcommand{\R}{\mathsf R}
\newcommand{\Good}{\mathsf{Good}}
\newcommand{\Case}{\mathsf{Case}}
\newcommand{\Inv}{\mathsf{Inv}}
\newcommand{\ev}{\operatorname{ev}}
\newcommand{\Add}{\mathsf{Add}}
\newcommand{\Mul}{\mathsf{Mul}}
\newcommand{\z}{\mathsf z}
\newcommand{\suc}{\mathsf s}
\newcommand{\nc}{\mathop{\neg}\nolimits_c}
\newcommand{\ac}{\mathbin{\wedge_c}}
\newcommand{\oc}{\mathbin{\vee_c}}
\newcommand{\iffc}{\mathbin{\leftrightarrow_c}}
\newcommand{\Up}{\mathsf{Up}}
\newcommand{\Da}{D_{\!a}}
\newcommand{\za}{\mathsf z_a}
\newcommand{\sa}{\mathsf s_a}
\newcommand{\adda}{\mathsf a_a}
\newcommand{\ma}{\mathsf m_a}
\newcommand{\Ga}{G_a}
\newcommand{\Ea}{E_a}
\newcommand{\Oa}{O_a}
\newcommand{\Gooda}{\mathsf{Good}_a}
\newcommand{\Casea}{\mathsf{Case}^{a}}
\newcommand{\Inva}{\mathsf{Inv}^{a}}
\newcommand{\nua}[1]{\nu^{a}_{#1}}

\usepackage{orcidlink}


% Keep the original article's body typography and all mathematics packages.
% Only title/author layout is delegated to the MLTT-S metadata class.
\usepackage{titling,authblk}
% Restore amsart's administrative-footnote order after titling's title.
\makeatletter
\@ifclassloaded{amsproc}{%
  % This is a manuscript date, not a journal's editorial receipt date.
  \renewcommand{\datename}{\textit{Date}:}%
  \apptocmd{\maketitle}{%
    \begingroup
    \let\thefootnote\@empty
    \let\markdownclass@dateprefix\@empty
    \ifx\@date\@empty\else
      \footnotetext{\@setdate}%
    \fi
    \ifx\@subjclass\@empty\else
      \footnotetext{\@setsubjclass}%
    \fi
    \endgroup
  }{}{\PackageError{paper}{Cannot display the AMS date/classification footnotes}{}}%
}{}
\makeatother
\setlength{\droptitle}{-1em}
\predate{\begin{center}\normalsize}
\makeatletter
\@ifclassloaded{amsproc}{%
  % Retain AMS body margins, but avoid the oversized first-page title gap.
  \setlength{\droptitle}{-2.5em}%
  % Avoid stacking titling's center-list and end-of-title skips on top of
  % the AMS abstract's own 20pt baseline-based separation.
  \postauthor{\end{tabular}\par\end{center}\removelastskip}%
  % amsproc stores the date but does not print it in its native title routine.
  \patchcmd{\@maketitle}
    {{\@bspredate \@date \@bspostdate}}{}{}{%
      \PackageError{paper}{Cannot suppress the non-AMS title date}{}%
    }%
  \patchcmd{\@maketitle}{\vskip 1.5em}{}{}{%
    \PackageError{paper}{Cannot remove extra titling space before AMS abstract}{}%
  }%
  % The metadata class sets this after loading the external preamble.
  \AtBeginDocument{\markdownclass@dateskip\z@}%
}{}
\makeatother
\usepackage[
  smartHybrid, fencedEnvs, headerAttributes, rawAttribute,
  inlineCodeAttributes, fencedCode, inlineNotes,
  citations, relativeReferences,
  % Avoid underscore-emphasis scanning across several native math subscripts.
  % Prose emphasis uses *...*, as in the manuscript.
  underscores=false,
  texMathDollars, texMathSingleBackslash
]{markdown}
